@@M0@@
Moreover, the additivity of @@M1@@ implies that @@M2@@, and the relations of the corollary are satisfied.

If @@M3@@ is the @@M4@@-algebra generated by generators @@M5@@ subject to the relations:
@@M6@@
there is therefore a surjective homomorphism from @@M7@@ onto @@M8@@. By 4.4, this homomorphism sends a basis of @@M9@@ to a basis of @@M10@@, so it is an isomorphism, which completes the proof of the corollary.

Let us treat the general case of a flag bundle of type @@M11@@. The usual flag bundle @@M12@@ can be obtained from @@M13@@ by successively constructing the flag bundles of the @@M14@@ (4.2). It follows that @@M15@@, considered as a module over @@M16@@, has a basis formed by the elements considered in 4.3.i), the @@M17@@ being the classes of the sheaves analogous to the @@M18@@ occurring in the splitting of the inverse images of the @@M19@@ on @@M20@@ (notation of 4.5 and 4.7). Thus @@M21@@ is a subring of @@M22@@; moreover, the additivity of @@M23@@ implies that:
@@M24@@
and @@M25@@ contains the sub-@@M26@@-algebra of @@M27@@ generated by the @@M28@@. Let @@M29@@ be this subalgebra; by 4.3.i), and what has just been said, there is in @@M30@@ a family of elements that forms a basis both for the @@M31@@-module structure and for the @@M32@@-module structure; consequently, @@M33@@, which shows that @@M34@@ is generated by the @@M35@@. The ideal of relations among the @@M36@@ is then given by 4.3.ii), taking
